La Tabla~\ref{tab:campos_wehoop} reúne una selección de campos tal como se identifican en los diccionarios publicados por SportsDataverse \cite{wehoopPbp,wehoopShots,wehoopLineups,wehoopPlayerSeason}. Se incluyen sus tipos de origen, conservando la distinción entre identificadores, características numéricas e información contextual.

{\small
\renewcommand{\arraystretch}{1.14}
\begin{longtable}{@{}p{3.2cm}p{3.2cm}p{1.65cm}p{4.9cm}@{}}
\caption{Variables originales documentadas en los conjuntos de datos.}\label{tab:campos_wehoop}\\
\toprule
\textbf{Variable} & \textbf{Origen} & \textbf{Tipo} & \textbf{Descripción}\\
\midrule
\endfirsthead
\multicolumn{4}{l}{\small Continuación de la Tabla~\ref{tab:campos_wehoop}}\\
\toprule
\textbf{Variable} & \textbf{Origen} & \textbf{Tipo} & \textbf{Descripción}\\
\midrule
\endhead
\midrule\multicolumn{4}{r}{\small Continúa en la siguiente página}\\\endfoot
\bottomrule\endlastfoot
\texttt{game\_id} & \texttt{pbp} & String & Clave que identifica cada encuentro.\\
\texttt{season} & \texttt{pbp} & Int64 & Año de la temporada.\\
\texttt{period} & \texttt{pbp} & Int64 & Periodo del evento, incluidos tiempos extra.\\
\texttt{clock} & \texttt{pbp} & String & Tiempo restante según el reloj del periodo.\\
\texttt{action\_number} & \texttt{pbp} & Int64 & Posición de la acción en el orden del partido.\\
\texttt{team\_id} & \texttt{pbp} & Int64 & Identificador del equipo al que se asocia la acción.\\
\texttt{person\_id} & \texttt{pbp} & Int64 & Persona vinculada al evento; no indica necesariamente a una defensora.\\
\texttt{action\_type} & \texttt{pbp} & String & Clasificación general del evento registrado.\\
\texttt{score\_home} & \texttt{pbp} & String & Marcador del equipo local después de la acción.\\
\texttt{score\_away} & \texttt{pbp} & String & Marcador visitante después de la acción.\\
\texttt{x\_legacy} & \texttt{shots} & Int64 & Coordenada horizontal del lanzamiento en el sistema de origen.\\
\texttt{y\_legacy} & \texttt{shots} & Int64 & Coordenada vertical del lanzamiento en el sistema de origen.\\
\texttt{shot\_result} & \texttt{shots} & String & Indica si el lanzamiento se convirtió o falló.\\
\texttt{player\_id} & \texttt{player\_season\_stats} & Int64 & Identificador persistente de jugadora.\\
\texttt{age} & \texttt{player\_season\_stats} & Float64 & Edad registrada para el periodo.\\
\texttt{ts\_pct} & \texttt{player\_season\_stats} & Float64 & Porcentaje de tiro verdadero.\\
\texttt{usg\_pct} & \texttt{player\_season\_stats} & Float64 & Porcentaje de uso ofensivo.\\
\texttt{off\_rating} & \texttt{lineups} & Float64 & Indicador ofensivo agregado de una alineación.\\
\texttt{def\_rating} & \texttt{lineups} & Float64 & Indicador defensivo agregado de una alineación.\\
\texttt{net\_rating} & \texttt{lineups} & Float64 & Diferencial de indicadores ofensivo y defensivo.\\
\texttt{poss} & \texttt{lineups} & Int64 & Posesiones asociadas al registro agregado.\\
\end{longtable}
}

\noindent\textbf{Descripción de la tabla.}
La Tabla~\ref{tab:campos_wehoop} constituye un diccionario resumido de variables de entrada, no una tabla de resultados del modelo. Los campos del registro de eventos permiten ordenar acciones y reconstruir el marcador, mientras que las estadísticas de temporada describen el perfil general de una jugadora. Por su parte, los ratings de \texttt{lineups} son agregados y no permiten inferir automáticamente el rendimiento de un quinteto ante cinco rivales específicos. La posición de los tiros no permite identificar trayectorias continuas de todas las jugadoras ni marcajes defensivos.